\section{The \texttt{MCPL\_input} McXtrace Component}
Source-like component that reads photon state parameters from an mcpl-file.

\subsection*{Identification}
\begin{itemize}
  \item \textbf{Author:} Erik B Knudsen
  \item \textbf{Origin:} DTU Physics
  \item \textbf{Date:} Aug 2016
\end{itemize}

\subsection*{Description}
Source-like component that reads photon state parameters from a binary mcpl-file.

MCPL is short for Monte Carlo Particle List, and is a new format for sharing events between e.g. MCNP(X), Geant4 and McXtrace .

When used with MPI, the --ncount given on the commandline is overwritten by \#MPI nodes x \#events in the file.

Example: MCPL\_input(filename=voutput,verbose=1,repeat\_count=1,E\_smear=0.1,pos\_smear=0.001,dir\_smear=0.01)

\subsection*{Input parameters}
Parameters in \textbf{boldface} are required; the others are optional.

\begin{longtable}{p{0.22\textwidth}p{0.12\textwidth}p{0.46\textwidth}p{0.14\textwidth}}
\toprule
\textbf{Name} & \textbf{Unit} & \textbf{Description} & \textbf{Default} \\
\midrule
\endhead
filename & str & Name of photon mcpl file to read & 0 \\
polarisationuse &   & If !=0 read polarisation vectors from file & 1 \\
verbose &   & Print debugging information for first 10 particles read & 1 \\
Emin & keV & Lower energy bound. Particles found in the MCPL-file below the limit are skipped & 0 \\
Emax & keV & Upper energy bound. Particles found in the MCPL-file above the limit are skipped & FLT\_MAX \\
repeat\_count & 1 & Repeat contents of the MCPL file this number of times. NB: When running MPI, repeating is implicit and is taken into account by integer division. MUST be combined sith the \_smear options! & 1 \\
E\_smear & 1 & When repeating events, make a Gaussian MC choice within E\_smear*E around particle energy E & 0 \\
pos\_smear & m & When repeating events, make a flat MC choice of position within pos\_smear around particle starting position & 0 \\
dir\_smear & deg & When repeating events, make a Gaussian MC choice of direction within dir\_smear around particle direction & 0 \\
preload &   & Load particles during INITIALIZE. On GPU preload is forced & 0 \\
\bottomrule
\end{longtable}

\subsection*{Links}
\begin{itemize}
  \item Component source code found in file \texttt{MCPL\_input.comp}.
\end{itemize}
\IfFileExists{misc/MCPL_input_static.tex}{\input{misc/MCPL_input_static.tex}}{}